\subsection{搜索}
\subsubsection{深度优先搜索}
DFS沿一个分支深入，结束后返回，继续尝试其他分支。状态要包含决定后续选择的信息；递归前修改的共享状态，返回后恢复。
下面用path保存已选前缀，used标记哪些数正在这个前缀中，依次选择1到n的未用数字。
% snippet:dfs
\begin{lstlisting}
// n <= 10; path[10], used[11] are global, initially zero.
void dfs(int depth) {
    if (depth == n) { visit(path, n); return; }
    for (int x = 1; x <= n; x++) {
        if (used[x]) continue;
        path[depth] = x; used[x] = 1;
        dfs(depth+1);
        used[x] = 0;
    }
}
\end{lstlisting}
调用dfs(0)开始。每个分支的有效路径长度是depth，下一次赋值会覆盖path[depth]，不必清空；used必须撤销，否则会阻止其他分支选择x。共有$n!$个完整排列，若visit遍历n项，总时间$O(n\cdot n!)$，数组与递归栈空间$O(n)$。
\subsubsection{基础剪枝}
剪枝排除不可能产生所需答案的分支，不能只凭当前结果“不好看”就停止。例如正整数子集求和到target时，已选和sum超过target后不可能恢复；sum加上所有剩余数仍小于target时也不可能达到目标。
\begin{lstlisting}
// rest is the sum of all still available positive numbers.
if (sum > target || sum + rest < target) return;
\end{lstlisting}
若允许负数，第一个条件不再可靠；每条剪枝都要由剩余选择能产生的范围推出。剪枝减少实际搜索量，一般不改变最坏情况的指数级复杂度。
\subsubsection{广度优先搜索}
BFS按距离从近到远逐层扩展，适用于每条边代价相同的最短路。用队列保存已发现但未扩展的顶点，首次入队时记录距离，避免重复入队。
% snippet:bfs
\begin{lstlisting}
// adj[u] stores neighbors; dist and q each have at least n slots.
for (int i = 0; i < n; i++) dist[i] = -1;
int head = 0, tail = 0;
dist[start] = 0; q[tail++] = start;
while (head < tail) {
    int u = q[head++];
    for (int i = 0; i < (int)adj[u].size(); i++) {
        int v = adj[u][i];
        if (dist[v] != -1) continue;
        dist[v] = dist[u] + 1;
        q[tail++] = v;
    }
}
\end{lstlisting}
顶点编号0到$n-1$，每个顶点最多入队一次。队首距离不大于后方顶点，第一次到达v时已经使用最少边数。不可达顶点保留$-1$。邻接表有n个顶点、m条边时，时间$O(n+m)$，总空间$O(n+m)$；边权不同时不能用普通BFS求最小费用。
